\documentclass[11pt,a4paper]{article}
\usepackage{amsmath,amssymb,amsthm}
\usepackage[margin=2.5cm]{geometry}
\usepackage{hyperref}

\newtheorem{definition}{Definition}[section]
\newtheorem{theorem}[definition]{Theorem}
\newtheorem{proposition}[definition]{Proposition}
\newtheorem{remark}[definition]{Remark}
\newtheorem{conjecture}[definition]{Conjecture}

\title{Hyperseed Formalization:\\
  Goertzel vs Epstein}
\author{ProtomegaTron (automated formalization)}
\date{2026-09-18}

\begin{document}
\maketitle

\begin{abstract}
We formalize the intellectual structure of Ben Goertzel's Substack article
``Goertzel vs Epstein'' (2026-02-20) within the Hyperseed ontology. The
article is a personal clarification of a funding relationship. Key mappings:
provenance documentation as fiber-base evidence trail, due diligence failure
as base verification failure, transparency as provenance exposure, and
funding as base resource with non-transmissive moral properties.
\end{abstract}

\tableofcontents

\section{Provenance as fiber-base evidence trail}
\label{sec:provenance}

\begin{theorem}[Provenance documentation --- claims 1--10, 21]
\label{thm:provenance}
The article is itself an exercise in provenance documentation: tracing the
evidence trail of a funding relationship between an AGI researcher and a
controversial funder.

In Hyperseed terms, the fiber (AGI research output --- OpenCog, Hyperon,
papers, projects) has a provenance history that includes the base (funding
sources). The funding base is a necessary resource for the fiber:
\[
  F_{\text{research}} \xrightarrow{\text{requires}} B_{\text{funding}}
\]

The specific funding relationship documented: approximately \$360K over
seventeen years of intermittent contact, from one of hundreds of funding
sources. The relationship was purely professional (funder-fundee), with
no social, personal, or criminal dimensions.

The key provenance fact: the fiber (research) was built on a base (funding)
that included a morally compromised source. The article documents this
provenance fully and transparently.
\end{theorem}

\section{Base verification failure}
\label{sec:verification}

\begin{proposition}[Due diligence --- claims 11--13, 23]
\label{prop:diligence}
Failing to verify the base (funding source) before building fiber (research)
on it is an epistemic integrity failure --- analogous to building inference
on unverified premises.

The hallucination bound from evidence conservation applies by analogy:
conclusions built on unverified premises inherit the uncertainty of those
premises. In the funding case:
\[
  \text{integrity}(\text{fiber}) \leq \text{verification}(\text{base})
\]

If the base is unverified ($\text{verification} < 1$), the fiber built
on it carries an integrity deficit --- not because the research itself
is flawed, but because the provenance chain includes an unverified link.

The due diligence failure was not unique: many scientists had similar
interactions with the same funder, who specifically targeted cutting-edge
science funding as a strategy. The systematic nature of the base
corruption (a funder deliberately targeting multiple researchers) means
many fiber bundles share the same compromised base link.
\end{proposition}

\section{Transparency as provenance exposure}
\label{sec:transparency}

\begin{theorem}[Provenance integrity --- claims 14--17, 22, 24--25]
\label{thm:integrity}
The fiber's provenance can't be retroactively cleaned or deleted --- it
must be honestly documented. Attempting to hide provenance is fiber
provenance deletion, which corrupts the integrity of the bundle.

Choosing transparency is choosing to expose the full provenance fiber:
\[
  \text{Transparency} = \text{expose}(\text{provenance}(F, B))
  \qquad \text{vs} \qquad
  \text{Concealment} = \text{delete}(\text{provenance}(F, B))
\]

This is the epistemic integrity equivalent of the anti-hallucination
theorems: honest inference requires visible provenance. An inference
engine that hides its evidence sources is less trustworthy than one
that exposes them, regardless of whether the evidence itself is good.

The moral quality of the base doesn't automatically corrupt the fiber.
Research funded by a morally compromised source is not itself morally
compromised --- the fiber has its own integrity independent of the
base's moral status. But the provenance connection remains and must
be acknowledged. This is a key distinction:
\[
  \text{moral}(\text{fiber}) \neq f(\text{moral}(\text{base}))
  \qquad \text{but} \qquad
  \text{provenance}(\text{fiber}, \text{base}) = \text{permanent}
\]
\end{theorem}

\section{Cross-references}

\begin{remark}[Connections to other formalizations]
\begin{itemize}
\item \textbf{Evidence Is to Logic What Energy Is to Physics} (2026-03-10):
  The hallucination bound and evidence conservation principles apply
  by analogy to funding provenance.
\item \textbf{LeCun's SAI Is a Special Case of AGI} (2026-03-07):
  Credit erasure as fiber provenance deletion --- the same principle
  applies to both funding provenance and intellectual provenance.
\item \textbf{Where Energy and AI Wars Collide} (2026-03-03):
  Funding as base resource, analogous to energy as base resource.
\end{itemize}
\end{remark}

\end{document}
